\begin{abstract}
Transparent objects break the brightness-constancy assumption that underlies
most segmentation pipelines: the appearance of glass is borrowed from the
background seen through it. A widely held intuition holds that \emph{physics}---
the refraction signature glass leaves between consecutive frames, and its
characteristic double-edge boundary response---should provide a cue that opaque
surfaces cannot. This paper is a controlled study of that intuition, not a
method that advocates it. Within a single architecture, \spectra, we add two
physics-motivated modules to a self-supervised \dinov{} ViT-S/14 backbone---an
Optical-Flow Consistency Violation (\ofcv) gate from frozen RAFT flow and a
non-trainable Boundary Resonance Field (\brf) Gabor prior---and isolate their
effect under controlled conditions. Our central finding is negative and, we
argue, useful: in a five-variant $\times$ 30-epoch $\times$ 3-seed ablation the
physics modules add no \emph{practically meaningful} accuracy---a paired
bootstrap bounds the largest physics effect at $+0.003$ \iou{} within a
$\pm0.01$ equivalence margin (TOST), with the Gabor prior statistically
indistinguishable from zero---and a real-video control shows the learned flow
cue is \emph{inert} even under genuine inter-frame parallax (the
\ofcv{} map moves by $<0.05\%$ and the prediction by $<0.001\%$ when the true
next frame replaces a duplicated one). The accuracy that does materialise is a
\emph{representation} effect: a \dinov{} backbone with a plain decoder---the
30-epoch ``neither'' ablation variant---reaches $0.932$ \iou{} on Trans10K
(versus our 10-epoch \spectra{} configuration at $0.9217$), and the proposed
modules leave this unchanged. We are explicit that our matched-budget comparison to
ImageNet-pretrained CNN baselines (\spectra{} $0.9217$ vs.\ DeepLabV3+ $0.8856$
binary foreground \iou) conflates backbone and pretraining corpus with method,
and we report it only as context for the operating point---never as a
contribution; the ablation, not this comparison, is the evidence. We position
\spectra{} as a refutation-style analysis: it quantifies, isolates, and explains
a popular but unverified inductive bias, and cautions against crediting appealing
physical priors when a strong self-supervised representation is the true cause.
We additionally provide zero-shot cross-dataset transfer (ClearPose, VGSD),
robustness over eight corruptions $\times$ five severities, an automated failure
taxonomy over all $4{,}428$ test images, and an interpretability analysis that
dissociates interpretable cues from accuracy.
\end{abstract}
